\documentclass[11pt,a4paper]{article}
\usepackage[margin=24mm]{geometry}
\usepackage{amsmath,amssymb,amsthm,array}
\setlength{\emergencystretch}{2em}
\newtheorem{lemma}{Lemma}
\newtheorem{theorem}[lemma]{Theorem}
\newcommand{\gen}[1]{\langle#1\rangle}
\title{The orbit budget of a small singleton prefix}
\author{Human Compression project: L044 continuation}
\date{}
\begin{document}\maketitle
Retain the complete critical singleton hypotheses and prefix-growth
proof. In particular $G\le A_n$ is transitive with connected string
C-marking, $G_0$ is the actual natural alternating group off $a$,
and all maximal ambient omissions are intransitive.
For $4\le g<r$ put
\[
 P=G_{<g},\quad Y=\operatorname{supp}(P),\quad
 C=a^P,\quad H=G_{\{2,\ldots,g-1\}},\quad b=a^{\rho_0}.
\]
Assume $P$ is intransitive on $Y$, $|Y|\le2g+1$, and each
colour $2,\ldots,g-1$ has at least two ambient-eligible edges.
These are the actual hypotheses at the first shared index of the
least later unique cut. All eligibility below refers to the full
ambient omission, never to an orbit projection.

\begin{lemma}\label{pair}
Suppose $\rho_0=(a\ b)(u\ v)$, the set $E=\{u,v\}$
is preserved by $\rho_1$ and $\rho_2$, and $\rho_1$ fixes
$a$ but moves $b$ to a point $x\notin\{a,b,u,v\}$.
Then $G$ is intransitive.
\end{lemma}
\begin{proof}
The $\gen{\rho_0,\rho_1}$-orbit of $a$ is $\{a,b,x\}$.
The set $E$ is an orbit of size two. Every other nontrivial orbit
has size two with $\rho_0$ acting trivially, since $\rho_0$
fixes every point outside $\{a,b,u,v\}$.
Any permutation centralizing both marks maps their orbits
equivariantly. It must preserve $E$, the only two-point orbit
on which $\rho_0$ acts nontrivially. Every mark of index at least
three centralizes both marks; the first three preserve $E$ by
hypothesis. Thus $G$ preserves the proper nonempty set $E$.
\end{proof}

\begin{theorem}\label{budget}
Let $f=|Y\cap\operatorname{Fix}(H)|$, and let $c$ count
the nontrivial $H$-orbits. Then $c\ge2$, and precisely the
following numerical possibilities remain:
\[
\begin{array}{c|c|c}
f&c&|Y|\\\hline
2&2&2g\text{ or }2g+1\\
3&2&2g+1\\
2&3&2g+1
\end{array}                                                     \tag{1}
\]
Moreover $\rho_0$ cannot exchange the points of any two-point
$H$-orbit. On the $H$-fixed points in $Y$, its action is exactly
$(a\ b)$. Its odd action on the remaining points consequently
has one of the following contributions: a free involution on a
preserved $H$-orbit of size $2q$, where $q\ge3$ is odd; or an
exchange of two $H$-orbits of the same odd size at least three.
This is a reduction to those possibilities, not their exclusion.
\end{theorem}
\begin{proof}
The group $H$ fixes $a,b$: all its marks fix $a$ and commute
with $\rho_0$. Both points lie in $Y$, so $f\ge2$.
The full ambient-interval theorem applies to $H$ and the given
eligible pairs. Every moved point of $H$ lies in $Y$, hence
\[
          2(g-2)\le |Y|-f-c,\qquad f+c\le5.              \tag{2}
\]
Also $c\ge1$ because the interval is nonempty and its marks
are nonidentity. Strict prefix growth gives $|C|\ge g+1$.

Suppose first that $c=1$, and denote the sole nontrivial $H$-orbit
by $U$. Then $f\le4$. Since $|C|\ge5>f$, the orbit $C$
meets $U$, hence contains it. A nontrivial $P$-orbit outside
$C$ must lie among the $H$-fixed points. Such an orbit has at
least two points, and $a,b$ already lie in $C$. Thus $f=4$
and the other $P$-orbit is $E=\{u,v\}$; the fixed set in
$Y$ is exactly $\{a,b,u,v\}$. Inequality (2) now forces
$|Y|=2g+1$ and $|U|=2g-3$.

The mark $\rho_0$ centralizes $H$ and preserves its unique
nontrivial orbit $U$. A centralizing involution on an odd transitive
orbit fixes that orbit pointwise: it has a fixed point, and commutation
then propagates fixation to the whole orbit. On the four remaining
points $\rho_0$ swaps $a,b$, so its evenness forces
$\rho_0=(a\ b)(u\ v)$. The mark $\rho_1$ preserves $E$
because $E$ is a $P$-orbit. Strict growth makes it move $b$;
its image is in $U$, since it fixes $a$ and preserves $E$.
Finally $\rho_2\in H$ fixes $E$ pointwise. Lemma~\ref{pair}
contradicts ambient transitivity. Thus $c\ge2$.

Suppose next that $E=\{u,v\}$ is an $H$-orbit on which
$\rho_0$ interchanges the two points. No edge of any colour
$2,\ldots,g-1$ on $E$ is ambient-eligible: its endpoints
are already joined by $\rho_0$ in the corresponding omission.
All eligible edges of that interval are therefore confined to
the union $V$ of the other nontrivial $H$-orbits.
The complete confinement lemma, including its actual orbit count,
gives
\[
  2(g-2)\le |V|-(c-1)=|Y|-f-c-1,
                      \qquad f+c\le4.                  \tag{3}
\]
There is at least one such other orbit, since every colour has
eligible edges. Since $f\ge2,c\ge2$, equality forces
$f=c=2$, $|Y|=2g+1$, and $V$ is one orbit of odd size
$2g-3\ge5$. The mark $\rho_0$ preserves $V$, because
its size differs from two, and fixes it pointwise by the same
centralizer argument. Thus again $\rho_0=(a\ b)(u\ v)$.
The orbit $C$ has at least five points, whereas $Y\setminus V$
has only four, so $V\subseteq C$. If $C$ met $E$ it would
contain all of $Y$, contrary to intransitivity of $P$ on $Y$.
Consequently $E$ is the other $P$-orbit, preserved by
$\rho_1$ and $\rho_2$. Strict growth makes $b^{\rho_1}$
lie in $V$. Lemma~\ref{pair} gives the contradiction again.

The table follows immediately from (2), $f\ge2,c\ge2$
and $|Y|\le2g+1$. Centralization of $H$ makes $\rho_0$
preserve its fixed set in $Y$. Its size is two or three, and
$\rho_0$ already swaps $a,b$, so it fixes any third point.
It therefore acts oddly on the nontrivial $H$-orbits.
On a preserved transitive orbit, a nonidentity centralizing
involution is free; its sign is odd exactly when the orbit has
size $2q$ with $q$ odd. The two-point case was just excluded,
so $q\ge3$. An exchange of two orbits contributes the sign
$(-1)^m$, where $m$ is their common size. An odd contribution
of that form thus has odd $m\ge3$. These possibilities exhaust
the odd sign and prove the asserted reduction.
\end{proof}

\paragraph{Attribution and boundary.}
The complete ambient-interval bound, confinement lemma and actual
singleton prefix growth are acquired inputs. The application to an
initial interval, its small fixed-set budget, and the uniform invariant-pair
exclusions are the present continuation. The source setting is
Cameron--Fernandes--Leemans--Mixer, \emph{Highest rank of a polytope
for $A_n$}, Propositions~5.28--5.32.
The two- and three-orbit cases in (1) have not been eliminated.
No faithful projected marking, projected IP, or pure alternating
prefix kernel is assumed. There is no new finite computation,
formal declaration, whole-original recovery or question credit.
\end{document}
